% ====================================================
% ====================================================
% USEPACKAGES
% ====================================================
% ====================================================

\usepackage[T1]{fontenc}
\usepackage[utf8]{inputenc}
\usepackage[english]{babel}

% tables
\usepackage{tabularx}
\usepackage{colortbl}
\usepackage{multirow}
\usepackage{makecell}

% tikz and colors
\usepackage{tikz}
\usepackage{xcolor}
\usepackage{pgf-pie}
\usepackage{pgfplots}
\usepackage{pgfplotstable}
\usepackage{tikzsymbols}

\usetikzlibrary{calc}
\usetikzlibrary{trees}
\usetikzlibrary{patterns}
\usetikzlibrary{shadings}
\usetikzlibrary{positioning}
\usetikzlibrary{intersections}
\usetikzlibrary{decorations.pathreplacing}

\usetikzlibrary{arrows}
\usetikzlibrary{arrows.meta}

\usetikzlibrary{shapes}
\usetikzlibrary{shapes.arrows}
\usetikzlibrary{shapes.callouts}
\usetikzlibrary{shapes.symbols}
\usetikzlibrary{shapes.geometric}

\usepgfplotslibrary{patchplots}
\usepgfplotslibrary{fillbetween}

% boxes
\usepackage[many]{tcolorbox}

% math packages and fonts
\usepackage{bm}
\usepackage{ccfonts}
\usepackage{eulervm}
\usepackage{amsmath}
\usepackage{amsfonts}
\usepackage{amssymb}
\usepackage{amsthm}
\usepackage{mathtools}
\usepackage{nicefrac}
\usepackage{slashed}
\usepackage{bbold}
\usepackage{array}
\usepackage{cancel}

% algorithms and listings
\usepackage[ruled,vlined,linesnumbered]{algorithm2e}
\usepackage{listings}
\usepackage{setspace}

\tcbuselibrary{listings}
\tcbuselibrary{breakable}
\tcbuselibrary{skins}

% misc
\usepackage{soul}
\usepackage{pifont}
\usepackage{skull}
\usepackage{multicol}
\usepackage{animate}
\usepackage{cleveref}
\usepackage{hyperref}
\usepackage{wasysym}
\usepackage[absolute,overlay]{textpos}
\usepackage[hang,flushmargin]{footmisc}
\usepackage[framemethod=tikz]{mdframed} 				% provides frame for framed figure and framed table in theme_2

% variable definitions
% =============================================
% variable deeptoc (deep or flat table of contents)
\newcommand{\dwDeepToc}[1]{
	\newboolean{deeptoc}
	\setboolean{deeptoc}{#1}
}

% figure and table name in caption
\addto\captionsenglish{%
  	\renewcommand{\figurename}{Abbildung}%
  	\renewcommand{\tablename}{Tabelle}%
}

% adjust font
% =============================================
\renewcommand*{\familydefault}{\sfdefault}

% framed Floatings
% =============================================
% define myBox
\newmdenv[
	innerlinewidth=0.05pt,
	roundcorner=4pt,
	linecolor=black,
	innerleftmargin=6pt,
	innerrightmargin=6pt,
	innertopmargin=6pt,
	innerbottommargin=6pt
]{mybox}

\newmdenv[
	innerlinewidth=0.05pt,
	roundcorner=4pt,
	linecolor=black,
	innerleftmargin=0pt,
	innerrightmargin=0pt,
	innertopmargin=0pt,
	innerbottommargin=-1pt
]{mytablebox}

% framed figure
\newcommand{\dwFigure}[3]{
	\begin{figure}
		\begin{mybox}
			\centering #1
  		\end{mybox}
  		\vspace{-4mm}
  		\caption{#2}
  		\label{#3}
	\end{figure}
	\vspace{-3mm}
}

% framed table
\newcommand{\dwTable}[4]{
	\begin{table}
		\begin{mytablebox}
			\renewcommand{\arraystretch}{#4} #1
		\end{mytablebox}
		\vspace{-2.5mm}
		\caption{#2}
		\label{#3}
	\end{table}
}

% centered X column in tabularx
\newcolumntype{Y}{>{\centering\arraybackslash}X}

% sections and subsections
% =============================================
% section command
\newcommand{\dwSection}[1]{
	\ifthenelse{\boolean{deeptoc}}{
		\section{#1}
	}{
		\section{#1}
		\subsection*{#1}
	}
}

% subsection command
\newcommand{\dwSubsection}[1]{\subsection{#1}}

% header
% =============================================
% green header
\newcommand{\dwHeader}[1]{%
	\begin{tcolorbox}[
		boxrule=0.2mm,
		boxsep=-1mm,
		lowerbox=ignored,
		colback=green,
		colframe=black,
		hbox
	]
		\textbf{#1}
	\end{tcolorbox}
}

% alert box
% =============================================
\newcommand{\dwAlertBox}[1]{
	\vspace*{2mm}\hspace*{0.25mm}
	\begin{minipage}[c]{0.05\textwidth}
		\begin{tikzpicture}
			\draw[thick,fill=lightgray] (4.2,0) -- (4.5,0.5) -- (4.8,0) -- (4.2,0) -- cycle;
			\fill (4.45,0.17) rectangle (4.55,0.37);
			\fill (4.55,0.08) arc (0:360:0.05);
		\end{tikzpicture}
	\end{minipage}
	\hfill
	\begin{minipage}[c]{0.92\textwidth}
		\begin{mybox}
			\textcolor{red}{\textbf{#1}}
		\end{mybox}
	\end{minipage}
}

% custom itemize environment
% =============================================
\let\tempone\itemize
\let\temptwo\enditemize
\renewenvironment{itemize}{\vspace*{1.5mm}\tempone\addtolength{\itemsep}{0.5\baselineskip}}{\temptwo}

% frames
% =============================================
\newenvironment{dwHeaderFrame}[1]{
	\subsubsection{#1}
	\begin{frame}{#1}
}{
	\end{frame}
}

% print title frame
% =============================================
\newcommand{\dwPrintTitle}{
	{\usebackgroundtemplate{%
		\tikz[overlay,remember picture] \node[opacity=0.9, at=(current page.center)] {
  			\includegraphics[height=\paperheight,width=\paperwidth]{../03_img/processor_red.jpg}
		};
	}
	\begin{frame}[plain]
		\begin{center}
			\begin{tcolorbox}[
				skin=enhanced,
	      			boxrule=0.6mm, boxsep=0mm,
				lowerbox=ignored,
				colback=orange!60!red, colframe=black,
				borderline={0.5pt}{3pt}{black}, borderline={1pt}{2pt}{red},
				width=\textwidth
			]
				\centering
				\Huge\textbf{\inserttitle}
			\end{tcolorbox}
			\vspace*{1.4cm}
			\begin{tcolorbox}[width=0.5\textwidth]
				\centering
				\textbf{\insertauthor} \\[2mm]
				\insertinstitute \\[2mm]
				\insertdate
			\end{tcolorbox}
			
			\begin{textblock}{1}(1,13.5)
				\includegraphics[scale=0.04]{../03_img/logo_dhbw}
			\end{textblock}
		\end{center}
	\end{frame}}
}

% print table of contents
% =============================================
\newcommand{\dwPrintToc}[1]{
	{\makeatletter
   		\setbeamertemplate{headline}[default]
   		\def\beamer@entrycode{\vspace*{-\headheight}}
	\makeatother

	\begin{frame}[allowframebreaks]
		\begin{tcolorbox}[
			skin=enhanced,
      			boxrule=0.6mm, boxsep=0mm,
			lowerbox=ignored,
			colback=orange!60!red, colframe=black,
			borderline={0.5pt}{3pt}{black}, borderline={1pt}{2pt}{red},
			width=\textwidth
		]
			\centering
			\huge\textbf{Inhaltsverzeichnis zur Vorlesung}
		\end{tcolorbox}
		\vspace{2mm}
		
		{\renewcommand{\baselinestretch}{1.4}
		\tableofcontents}
	\end{frame}}
}

% redefine cref (add a hyperlink)
\let\chyperref\cref % save original command under a new name
\renewcommand{\cref}[1]{\hyperlink{#1}{\textcolor{blue}{$\Rightarrow$ \chyperref{#1}}}}

\newcommand{\externalurl}[2]{\href{#1}{\textcolor{blue}{$\Rightarrow$ #2}}}




% ====================================================
% ====================================================
% OPTIONS
% ====================================================
% ====================================================

% select theme
\usetheme{mystyle}

% number of levels in toc
\dwDeepToc{false}

\newcommand{\modulkatalog}{\externalurl{https://www.dhbw.de/fileadmin/user/public/SP/MA/Wirtschaftsinformatik/Data_Science.pdf}{Modulkatalog}}
\newcommand{\sharelatex}{\externalurl{https://de.sharelatex.com/templates/journals/american-institute-of-physics-aip}{Vorlage American Institute of Physics (AIP)}}
\newcommand{\overleaf}{\externalurl{https://de.overleaf.com/latex/templates/tagged/report}{Report Vorlagen Overleaf}}


% ====================================================
% ====================================================
% PRESENTATION DATA
% ====================================================
% ====================================================

\title[Data Exploration Project]{Data Exploration Project}
\author{M.\,Sc. Daniel Wehner}
\date{Sommersemester 2020}
\institute{SAP\,SE / DHBW Mannheim}

% ====================================================
% ====================================================
% BEGIN OF DOCUMENT
% ====================================================
% ====================================================

\begin{document}

% Title frame
%______________________________________________________________________
\dwPrintTitle

% Agenda
%______________________________________________________________________
\dwPrintToc

% Section: Planung und Organisatorisches
%______________________________________________________________________
\dwSection{Planung und Organisatorisches}

% Grundsätzliches zum Projekt
\begin{dwHeaderFrame}{Grundsätzliches zum Projekt}
	\begin{itemize}
		\item Name der Veranstaltung: \textit{Data Exploration Project}
		\item Laut \modulkatalog{} beträgt der Workload des Projekts \textbf{pro Person}:
		\begin{itemize}
			\item Präsenzzeit: 27 Stunden
			\item Selbststudium: 47 Stunden
		\end{itemize}
		\item Definition des Projekts aus dem \modulkatalog:
		\begin{quote}
			\glqq{}Anwendung von Methoden und Verfahren des maschinellen Lernens auf eine vorgegebene Datenbasis unter Laborbedingungen. Verwendung von üblichen Repositorien wie Hadoop/Spark/Flink/Mahout, Python-RASBT, R, etc.
			\textbf{Ein besonderer Fokus soll auf einer ganzheitlichen wirtschaftsinformatischen Betrachtung liegen.} Es soll dabei neben der informatischen Betrachtung auch der betriebswirtschaftliche Nutzen, z.\,B. anhand eines Use Cases, betrachtet werden.\grqq{}
		\end{quote}
	\end{itemize}
\end{dwHeaderFrame}


% Bearbeitung und Abgabe
\begin{dwHeaderFrame}{Bearbeitung und Abgabe}
	\dwHeader{Bearbeitung}
	\begin{itemize}
		\item Das Projekt ist in Gruppen von \textbf{maximal vier Studierenden} zu bearbeiten (bei Bedarf können auch kleinere Gruppen gebildet werden, größere Gruppen sind jedoch nicht erlaubt).
		\item Jede Gruppe bearbeitet ein anderes Thema.
		\item Mögliche Projektthemen: siehe ...
		\item Bzgl. der zu verwendenden Technologien werden keine Einschränkungen gemacht.
	\end{itemize}
	
	\dwAlertBox{Bei der Themenwahl ist darauf zu achten, den Umfang weder zu gering noch zu groß zu wählen!}
\end{dwHeaderFrame}


\begin{frame}
	\dwHeader{Zwischenpräsentation}
	\begin{itemize}
		\item Die Präsentation beträgt \textbf{maximal 10 Minuten} und dient dem Zweck, die grobe Konzeption des Projektvorhabens darzulegen und zu präsentieren.
		\item Sie geht \textbf{nicht} mit in die Endwertung ein.
		\item Von Seiten der Projektgruppen sind zu diesem Termin keine Arbeitsergebnisse oder Dokumente einzureichen.
	\end{itemize}
	
	\vspace*{3mm}
	\dwHeader{Projektabschluss am Ende des Semesters}
	\begin{itemize}
		\item Am Ende des Semesters sind eine Abschlusspräsentation und ein Projektreport anzufertigen.
		\item Die Abschlusspräsentation beträgt \textbf{25 Minuten}.
		\item Der Projektreport ist gemäß den Regeln des wissenschaftlichen Arbeitens anzufertigen und weist einen Umfang von \textbf{minimal 3 und maximal 4 Seiten} (ohne Abbildungen) auf.
		\item Die Termine können der \cref{tab:ablauf-projekt} entnommen werden.
	\end{itemize}
\end{frame}


\begin{frame}
	\begin{itemize}
		\item \textbf{Tipp:} Online gibt es bereits diverse \LaTeX-Vorlagen für die Anfertigung von Projektreports; z.\,B.:
		\begin{itemize}
			\item ShareLaTeX: \sharelatex
			\item Overleaf: \overleaf
		\end{itemize}
		\item Der Projektreport deckt mindestens folgende Bestandteile ab:
		\begin{itemize}
			\item Thema und Motivation
			\item Related Work (\textit{welche wissenschaftlichen Publikationen gibt es zu diesem Thema?})
			\item Verwendete Technologien und Bibliotheken (z.\,B. \texttt{scikit-learn}, \texttt{tensorflow}, ...)
			\item Anmerkungen zum Quellcode (\textit{wie ist der Code auszuführen und was gibt es zu beachten?})
			\item Präsentation der Ergebnisse
			\item Kritische Bewertung der Ergebnisse (\textit{\glqq{}lessons learned\grqq{}})
		\end{itemize}
	\end{itemize}

	\dwAlertBox{Die komplette Abgabe besteht aus dem Projektreport (*.pdf), der Abschlusspräsentation (*.pdf), sowie dem im Rahmen des Projekts erstellten Quellcode.}
\end{frame}


% Anwesenheitspflicht
\begin{dwHeaderFrame}{Anwesenheitspflicht}
	\begin{itemize}
		\item Mit wenigen Ausnahmen ist die Anwesenheit der einzelnen Projektgruppen an den Terminen \textbf{freiwillig}.
		\begin{itemize}
			\item Es sollte jedoch regelmäßig Rücksprache bzgl. des Zwischenstands gehalten werden.
			\item Bitte melden Sie sich \textbf{rechtzeitig} und \textbf{eigenverantwortlich} (z.\,B. per E-Mail: \texttt{daniel.wehner@sap.com}), falls von Ihrer Seite aus Diskussionsbedarf besteht.
		\end{itemize}
		\item \textbf{Anwesenheitspflicht} besteht an folgenden Terminen (siehe \cref{tab:ablauf-projekt}):
		\begin{itemize}
			\item Zwischenpräsentation
			\item Finale Präsentation und Abgabe
		\end{itemize}
	\end{itemize}
\end{dwHeaderFrame}


% Bewertung
\begin{dwHeaderFrame}{Bewertung des Projekts}
	\begin{itemize}
		\item Nur rechtzeitig eingereichte Dokumente können bewertet werden.
		\item Die einzelnen Bestandteile werden folgendermaßen gewichtet:
		\begin{itemize}
			\item Quellcode und Ergebnisse (70\,\%)
			\item Projektreport (20\,\%)
			\item Abschlusspräsentation (10\,\%)
		\end{itemize}
	\end{itemize}

	\dwAlertBox{Wichtig: Das Fehlen einer Teilabgabe führt zu erheblichem Punkteabzug (unter Umständen auch zum Nichtbestehen der gesamten Veranstaltung)!}
\end{dwHeaderFrame}


% Zeitlicher Ablauf des Projekts
\begin{dwHeaderFrame}{Zeitlicher Ablauf des Projekts}
	\dwTable{
		\begin{tabularx}{\textwidth}{ c | Y | c }
			\textbf{Termin} 		&
			\textbf{Bemerkung} 	&								
			\textbf{Anwesenheitspflicht}									\\ \hline\hline
			xx.xx.2020			& 	\textbf{Einführung}				&	ja	\\ \hline
			xx.xx.2020			&	\textit{Bedarf bitte anmelden}		& 	nein	\\ \hline
			xx.xx.2020			& 	\textbf{Zwischenpräsentation (alle Gruppen)}
															&	ja	\\ \hline
			xx.xx.2020			&	\textit{Bedarf bitte anmelden}		&	nein	\\ \hline
			xx.xx.2020			&	\textit{Bedarf bitte anmelden}		&	nein	\\ \hline
			xx.xx.2020			&	\textit{Bedarf bitte anmelden}		&	nein	\\ \hline
			xx.xx.2020			&	\textit{Bedarf bitte anmelden}		&	nein	\\ \hline
			xx.xx.2020			& 	\textbf{Abschlusspräsentationen und Abgabe (alle Gruppen)}
															& 	ja
		\end{tabularx}
	}{Ablaufplan des Projekts}{tab:ablauf-projekt}{1.2}
\end{dwHeaderFrame}


% Section: Themen
%______________________________________________________________________
\dwSection{Themen}

% Eigene Themen
\begin{dwHeaderFrame}{Eigene Themen}
	\begin{itemize}
		\item Laut \modulkatalog{} soll der Fokus auf einer \textit{\glqq{}ganzheitlichen wirtschaftsinformatischen Betrachtung\grqq{}} liegen,
			und auch dem betriebswirtschaftlichen Aspekt Rechnung getragen werden.
		\item Es ist grundsätzlich erlaubt und auch erwünscht, \textbf{eigene Themenvorschläge} einzubringen.
		\item Eigene Themen müssen natürlich vorher genehmigt werden.
		\item Auf der nächsten Folie sind einige Projektvorschläge aufgelistet, falls einzelne Gruppen kein eigenes Thema finden sollten.
	\end{itemize}
\end{dwHeaderFrame}

% Themenvorschläge
\begin{dwHeaderFrame}{Themenvorschläge}
	\begin{itemize} 
		\item \texttt{Sentiment Analyse von Kundenrezessionen} (kommt das Produkt beim Kunden gut oder schlecht an?)
		\item \texttt{Vorhersage von Aktienkursen} (falls Sie reich werden möchten, ist das ein guter Anfang)
		\item \texttt{Recommender Systems} (z.\,B. \textit{Collaborative Filtering})
		\item \texttt{Baue deinen eigenen Chatbot} (z.\,B. zur automatischen Beantwortung von Kundenfragen)
		\item \texttt{Speech Recognition}
		\item \texttt{Analyse medizinischer Scans zur Krankheitsdiagnose}
	\end{itemize}
\end{dwHeaderFrame}

\end{document}